%% kapa: exam=1 round=26 year=2015 period=2 form=A subject=accounting subject_ko=회계학 seq=32 number=72
%% source: Q-Net 감정평가사 기출문제 (cst00309) artlSeq=5138001 file=extracted/1008/5138001/제26회_감정평가사 1차_기출/제26회_감정평가사 1차_기출/2교시_감평_A형.pdf
%% conversion: kordoc 4.24.51 -> tools/gen1.py (auto); review=proofread-sonnet; flags=none
\begin{question}{72}
(주)감평은 두 개의 보조부문(X부문, Y부문)과 두 개의 제조부문(A부문, B부문)으로 구성되어 있다. 각각의 부문에서 발생한 부문원가는 A부문 ￦100,000, B부문 ￦200,000, X부문 ￦140,000, Y부문 ￦200,000이다. 각 보조부문이 다른 부문에 제공한 용역은 다음과 같다.

{\small\centering\begin{tabular}{|c|c|c|c|c|}\hline
\multirow{2}{*}{\makecell[r]{사용부문\\[0.6em]제공부문}} & \multicolumn{2}{c|}{보조부문} & \multicolumn{2}{c|}{제조부문} \\ \cline{2-5}
 & X부문 & Y부문 & A부문 & B부문 \\ \hline
\makecell{X부문\\(kWh)} & - & 50,000 & 30,000 & 20,000 \\ \hline
\makecell{Y부문\\(기계시간)} & 200 & - & 300 & 500 \\ \hline
\end{tabular}\par}
\medskip
(주)감평이 단계배부법을 이용하여 보조부문원가를 제조부문에 배부할 경우, A부문과 B부문 각각의 부문원가 합계는? (단, 배부 순서는 Y부문의 원가를 먼저 배부한다.)

\noindent\hspace*{2em}\makebox[10em][l]{\uline{A부문원가 합계}}\uline{B부문원가 합계}\par
\begin{choices}
  \item \makebox[10em][l]{￦168,000}￦172,000
  \item \makebox[10em][l]{￦202,000}￦328,000
  \item \makebox[10em][l]{￦214,000}￦336,000
  \item \makebox[10em][l]{￦244,000}￦356,000
  \item \makebox[10em][l]{￦268,000}￦372,000
\end{choices}
\end{question}
\answer{5}
